\subsection{The proper abelian image is one global boundary class}

\begin{statement}{A.163}{Proved}
{A complete classification in the proper abelian cover quotient}
Use \(N_{\rm all}\), \(E_D\), and \(E_U\) from
A.160, and put \(J=\operatorname{Jac}(\overline D)\).
Take \(P_0\) as Abel--Jacobi origin and write
\(\iota(P)=[P-P_0]\).  There is an actual isomorphism
of arithmetic extensions
\[
                 E_D/N_{\rm all}\simeq TJ\rtimes G_\Q,
\tag{A.163.1}
\]
where the right side is split by \(\iota(P_0)=0\).

Every continuous section of \(E_U/N_{\rm all}\) in
the combined class \([64a^2]\) which locally matches
actual point or cusp sections is geometrically conjugate
to the image of an actual rational cusp section of the
tripod.  More precisely, choose any global fixed sheet
of the degree-eight cover, giving a lift to
\(E_D/N_{\rm all}\).  There is then one
\[
                         P\in\{P_0,P_1,P_{1/2},P_\infty\}
\tag{A.163.2}
\]
such that this lifted section is geometrically conjugate
to the proper point image at \(P\) in (A.163.1).
Every selected proper local centre at a finite place
is this same \(P\).  The resulting images are exactly
the images of these four rational boundary sections;
no distinctness assertion is made after the larger
\(\Pi/N_{\rm all}\)-conjugacy relation downstairs.
For a locally matching full metabelian section in this
class, this conclusion concerns its image under the
actual quotient map to \(E_U/N_{\rm all}\).

The two-primary arithmetic of the old elliptic factor
and of \(J\) can also be sharpened:
\[
 \begin{gathered}
 \operatorname{Sel}^{2}(E/\Q)=E(\Q)/2E(\Q)
                                  \simeq(\Z/2)^4,
       \qquad \Sha(E/\Q)[2^\infty]=0,\\
 \operatorname{rank}J(\Q)=2,
       \qquad 4\Sha(J/\Q)[2^\infty]=0,
       \qquad T_2\Sha(J/\Q)=0.
 \end{gathered}
\tag{A.163.3}
\]
Here \(E:y^2=x^3-1146^2x\) is A.143's rank-two
factor.  The middle two-primary group in (A.163.3)
is finite, without an asserted order.  No odd-primary
Tate--Shafarevich assertion for \(E\) or \(J\) is
used in the all-coefficient classification of the
particular finite-support classes in (A.163.1).
\end{statement}

\begin{proof}
\emph{The proper abelian extension and its actual Kummer map.}
The compactification quotient in A.157 is
\(H/\mathcal I=\pi_1(\overline D_{\overline\Q})\).
The same finite-cover correspondence gives the arithmetic
quotient: a cover with trivial boundary inertia extends
by normalization to a finite etale cover of
\(\overline D\), and conversely every such cover has
trivial boundary inertia.  Thus
\[
 E_D/N_{\rm all}
   =\pi_1(\overline D/\Q)\big/
       \overline{[\pi_1(\overline D_{\overline\Q}),
        \pi_1(\overline D_{\overline\Q})]}.
\tag{A.163.4}
\]
The Abel--Jacobi morphism \(\iota\) induces the
geometric isomorphism from the abelian group in this
extension to \(TJ\).  One literal verification uses
the complex period definition of \(J\): the map
\(P\mapsto\int_{P_0}^P\) identifies the integral
first homology of the curve with the period lattice
of its Jacobian.  Taking profinite completions and
the finite etale comparison gives the asserted map
on geometric fundamental groups.  Equivalently all
finite geometric abelian covers are the pullbacks
of the corresponding isogeny covers of \(J\),
as in the proof of Stoll's Theorem 6.4
\cite[Theorem 6.4 and its proof]{LegStoll2007}.
The multiplication covers are cofinal in these
isogeny covers and identify the target with \(TJ\).
All maps are defined over \(\Q\), and \(P_0\)
maps to zero.  They therefore identify the full
arithmetic extensions and their splittings, proving
(A.163.1).

Fix the global sheet specified in the statement.
A.160 proves that the section lifts into
\(E_D/N_{\rm all}\) and that its actual proper
local centres \(P_p\) lie in the four-point set
(A.163.2) for every finite prime.  Under (A.163.1)
write the section as a continuous cocycle
\(\eta\in H^1_{\rm cts}(\Q,TJ)\).  The actual
local fixed-sheet maps of A.160 identify its finite
components with
\(\kappa_m(\iota(P_p))\), and with the appropriate
real Kummer component at infinity.  Hence
\[
             (\eta_m)_m\in
              \widehat{\operatorname{Sel}}(\Q,J)
                 :=\varprojlim_m\operatorname{Sel}^m(J/\Q).
\tag{A.163.5}
\]
The local value of this inverse system is the actual
point \(\iota(P_p)\) in \(J(\Q_p)\): full local
Kummer injectivity is proved by the formal subgroup
argument of A.160, which applies in the same way to
an abelian variety.

Apply the opened primary Proposition 3.6 of Stoll
to the finite rational subscheme
\[
                  Z=\iota\{P_0,P_1,P_{1/2},P_\infty\}
                         \subset J.
\tag{A.163.6}
\]
That proposition says that an element of the full
inverse Selmer group whose actual local values lie
in a finite rational subscheme at density one is
itself the Kummer image of a point of that subscheme
\cite[Proposition 3.6, p.~9 of the author version]{LegStoll2007}.
It has no finiteness hypothesis on the Tate--Shafarevich
group.  Its closely related Proposition 3.7 states
injectivity of the inverse Selmer localization map
at density one.  We have verified its finite-support
hypothesis here at every finite place.  It follows
that \((\eta_m)_m=\kappa(\iota(P))\) for a point
\(P\) in (A.163.2).  The finite-zero-cochain
compactness argument of A.154 identifies the original
continuous classes as well:
\[
                              \eta=\kappa(\iota(P)).
\tag{A.163.7}
\]
This is equality of classes of sections under
conjugation by the actual geometric group \(TJ\)
in (A.163.1).

The Abel--Jacobi map is injective on geometric
points.  Indeed equality of \([P-P_0]\) and
\([Q-P_0]\) for distinct points would give a
rational function with divisor \(P-Q\), hence a
degree-one map of the genus-three curve to
\(\mathbf P^1\), which is impossible.  Thus local
full Kummer injectivity and (A.163.7) force
\(P_p=P\) at every finite place.  At the real place
the statement concerns the abelianized component,
as required by the Selmer convention.

Choose any full tangential section at the rational
cusp \(P\) of \(D^\circ\).  Compactification
kills its boundary inertia, so its image in
(A.163.1) is the class on the right of (A.163.7).
Its pushforward through the actual finite cover
\(D^\circ\to U\) is an actual cusp section at
the image prescribed in (A.160.21).  This proves
the stated geometric conjugacy in the quotient.
Conversely these cusp sections have full Selmer
localizations and a fixed sheet in the degree-eight
cover, so each supplies a section in the stated
domain.  The argument is carried out after choosing
one global sheet and makes no claim about whether
the four resulting images remain distinct after
additional deck conjugation downstairs.

\emph{The complete two-Selmer group of the old elliptic factor.}
Retain the actual dual isogenies
\(\phi:E\to F\) and \(\psi:F\to E\) in
(A.143.13), with \(F:v^2=u^3+573^2u\) and
\(\psi\phi=[2]\).  To fix orientation, write
\[
 \operatorname{Sel}^{\psi}(E/\Q)
                       \subset H^1(\Q,F[\psi]),\qquad
 \operatorname{Sel}^{\phi}(F/\Q)
                       \subset H^1(\Q,E[\phi]).
\tag{A.163.8}
\]
Their local conditions are, respectively, the
connecting images of \(E(\Q_v)\) and \(F(\Q_v)\).
The actual connecting maps are the \(x\)- and
\(u\)-squareclasses of A.143, including the special
values at the origin and at \((0,0)\).  Their
quartic representatives are
\(N^2=dU^4+(B/d)V^4\) for
\(B=-1146^2\) and \(B=573^2\), respectively.

The valuation calculation in A.143 applies to
each local connecting image, not only to a global
point.  Away from the prime support of \(B\),
the valuation of the connecting class is even;
the two special points have this property too.
Consequently
\[
 \operatorname{Sel}^{\psi}(E/\Q)
                  \subset\langle-1,2,3,191\rangle,
\tag{A.163.9}
\]
which has sixteen elements.  All sixteen already
occur from rational points by the four independent
classes of \((0,0),(1146,0),P,Q\) in A.143.
Thus (A.163.9) is equality.

On the \(F\)-side, its real local condition makes
the class positive.  The valuation calculation
leaves only \(1,3,191,573\).  The three exclusions
in A.143 are local exclusions and prove
\[
                         \operatorname{Sel}^{\phi}(F/\Q)=0.
\tag{A.163.10}
\]
For completeness, scale a putative local quartic
point so that \(U,V\) are integral with at least
one a unit.  At three, the classes three and 573
give a right side of valuation one, since the
parenthesis is respectively \(U^4+191^2V^4\) or
\(U^4+V^4\), nonzero modulo three.  At 191 the
remaining class gives \(191(U^4+9V^4)\).  Its
parenthesis cannot vanish for primitive coordinates,
because \(-9\) is a nonsquare modulo 191.
The right side again has valuation one.  These
checks include the coordinate-zero cases.

The exact finite-module sequence
\[
           0\longrightarrow E[\phi]\longrightarrow E[2]
                     \xrightarrow{\phi}F[\psi]\longrightarrow0
\tag{A.163.11}
\]
sends \(\operatorname{Sel}^2(E/\Q)\) into
\(\operatorname{Sel}^{\psi}(E/\Q)\).
Its kernel comes from
\(\operatorname{Sel}^{\phi}(F/\Q)\).  To check
the local condition in this last assertion, let
a kernel class lift to \(\alpha\in H^1(\Q,E[\phi])\).
At a place its two-Kummer point \(P_v\) has zero
\(\psi\)-connecting image, so \(P_v=\psi Q_v\)
for some \(Q_v\in F(\Q_v)\).  The image of
\(\kappa_\phi(Q_v)\) in \(H^1(\Q_v,E[2])\)
is its two-Kummer class.  The difference from
\(\alpha_v\) is a boundary coming from
\(F[\psi](\Q_v)\) in (A.163.11).  That boundary
is itself the \(\phi\)-Kummer image of this
torsion point.  Hence \(\alpha_v\) satisfies the
\(\phi\)-Selmer condition at every place.
Equation (A.163.10) proves the asserted injection.
It bounds \(\operatorname{Sel}^2(E/\Q)\) by
sixteen, and A.143 proved that \(E(\Q)/2E(\Q)\)
already has sixteen elements.  The finite Kummer
sequence therefore gives \(\Sha(E/\Q)[2]=0\).
Any nonzero two-primary torsion element has a
nonzero multiple of order two, so
\(\Sha(E/\Q)[2^\infty]=0\), as claimed.

\emph{An actual rational isogeny for the genus-three Jacobian.}
Over \(K=\Q(\sqrt2)\), the three quotient involutions
are
\[
 j_0(T,W)=(T,-W),\qquad
 j_\pm(T,W)=\left(\frac1{2T},
                            \frac{\pm W}{\sqrt2\,T}\right).
\tag{A.163.12}
\]
They commute, have order two, and satisfy
\(j_+j_-=j_0\).  Their quotients are exactly
\(q_{\rm old},q_+,q_-\).  The three differential
lines in A.150 have no common invariant.  This
can be checked directly on
\(dT/W^2,dT/W^3,TdT/W^3\): \(j_0\) acts by
\(1,-1,-1\), while \(j_+\) sends these forms to
\(-dT/W^2,-\sqrt2\,TdT/W^3,-dT/(\sqrt2 W^3)\).
The action of \(j_-\) is the product of those two
actions.  Their group sum is zero.  In characteristic
zero a homomorphism of abelian varieties with zero
differential is zero, so on \(J_K\)
\[
                         1+j_0+j_++j_-=0.
\tag{A.163.13}
\]
Define the actual divisor maps
\[
 \begin{aligned}
 \Phi:J&\longrightarrow E\times A,
                &\Phi_K&=(q_{{\rm old},*},q_{+,*},q_{-,*}),\\
 \Psi:E\times A&\longrightarrow J,
                &\Psi_K(R,S_+,S_-)
                         &=q_{\rm old}^*R+q_+^*S_++q_-^*S_-.
 \end{aligned}
\tag{A.163.14}
\]
Field conjugation exchanges the last two maps and
the two factors of \(A_K\); it fixes the first.
Thus both formulas descend to the displayed
\(\Q\)-morphisms.  Curve pullback and norm of
divisors give
\(q_j^*q_{j,*}=1+j\) and
\(q_{j,*}q_j^*=[2]\), including ramified fibres
with their multiplicities.  Equations (A.163.13)
and (A.163.14) yield
\[
                            \Psi\Phi=[2]_J,
                           \qquad\Phi\Psi=[2]_{E\times A}.
\tag{A.163.15}
\]
For the second equality, the off-diagonal maps
vanish: the image of \(q_k^*\) is fixed by \(j_k\)
and acted on by \(-1\) by either other involution,
as follows from (A.163.13).  The pushforward
\(q_{j,*}\) is invariant under \(j_j\); hence
\(2q_{j,*}q_k^*=0\) for \(j\ne k\).
The group of homomorphisms of abelian varieties
has no torsion, so this map is zero.
Both maps in (A.163.14) are consequently isogenies.
They prove \(\operatorname{rank}J(\Q)=2\), using
A.143 and A.160, without base-changing the
rank-two factor to a further number field.

The maps on \(\Sha\) preserve local triviality.
Their target \(\Sha(E\times A/\Q)[2^\infty]\)
is killed by two: its \(E\)-part is zero by
(A.163.3), and its \(A\)-part is A.154's
\((\Z/2)^2\) through the explicit Shapiro map.
Applying \(\Psi\Phi=[2]\) therefore gives
\(4\Sha(J/\Q)[2^\infty]=0\).
The fixed-level group \(\Sha(J/\Q)[4]\) is
finite by the Selmer finiteness argument of A.160,
so this two-primary subgroup is finite and has
zero Tate module.  This proves (A.163.3).
It also gives the exact compatible-two-Selmer
identification
\(\varprojlim_r\operatorname{Sel}^{2^r}(J/\Q)
=\widehat{J(\Q)}_2\).
The earlier all-coefficient equality (A.163.7)
instead used Stoll's finite-support proposition
for that specific class; it did not replace
\(\widehat{\operatorname{Sel}}(\Q,J)\) by the
full completed Mordell--Weil group.

The local quartic exclusions, the four independent
point squareclasses, and the exact involution and
differential identities are independently checked in
\path{code/harmonic_proper_abelian_certificate.py}
with retained output.
\end{proof}
